\section{Results}
\label{sec:results}

We present results for all four primary hypotheses, interaction effects, and supplementary analyses. All hypotheses were pre-registered on OSF prior to data collection; all reported $p$-values are two-tailed with Bonferroni-Holm correction unless noted.

\subsection{Main Effects: All Four Factors Significantly Predict Norm Emergence}
\label{sec:main-effects}

Table \ref{tab:main-effects} shows mean BCS by factor level. All four main effects were significant.

\begin{table*}[t]
\caption{Mean Behavioral Convergence Score (BCS) by Factor Level, with MANOVA Results}
\label{tab:main-effects}
\vskip 0.12in
\begin{center}
\begin{small}
\begin{tabular}{llcrrr}
\hline
\textbf{Factor} & \textbf{Level} & \textbf{Mean BCS} & \textbf{SD} & \textbf{$\eta^2$} & \textbf{$p$ (Holm)} \\
\hline
Feedback (F1) & HIGH & 0.52 & 0.19 & 0.14 & $<$0.001 \\
              & LOW  & 0.19 & 0.11 & —    & —              \\
\hline
Comm. Structure (F2) & HIGH & 0.68 & 0.22 & 0.31 & $<$0.001 \\
                      & LOW  & 0.09 & 0.07 & —    & —              \\
\hline
Task Framing (F3) & Cooperative & 0.44 & 0.21 & 0.06 & 0.003 \\
                  & Neutral     & 0.31 & 0.17 & —    & —              \\
\hline
Pop. Stability (F4) & Stable & 0.48 & 0.20 & 0.09 & $<$0.001 \\
                     & Churn  & 0.27 & 0.15 & —    & —              \\
\hline
\end{tabular}
\end{small}
\end{center}
\vskip -0.1in
\end{table*}

\textbf{H1 is confirmed}: Feedback (F1=HIGH) produced mean BCS = 0.52 (95\% CI [0.48, 0.56]) vs. 0.19 [0.15, 0.23] in LOW conditions ($t(158) = 14.2$, $p < 0.001$, Cohen's $d = 1.79$). However, no F1-only condition reached the 0.70 norm emergence threshold within 20 rounds, confirming that feedback loops are \emph{necessary but not sufficient} for norm crystallization.

\textbf{H2 is confirmed with the largest effect}: Communication structure (F2=HIGH) produced mean BCS = 0.68 [0.63, 0.73], significantly exceeding all other factors ($\eta^2 = 0.31$). Comparing F2 to F1: $\Delta\eta^2 = 0.17$, $F(1, 144) = 35.6$, $p < 0.001$. This establishes F2 as the dominant environmental determinant.

\textbf{H3 is confirmed}: Cooperative task framing (F3=HIGH) increased mean BCS by 0.13 [0.04, 0.22] relative to neutral framing ($t(158) = 2.99$, $p = 0.003$, $d = 0.47$). The effect was more pronounced when combined with F2=HIGH ($\Delta$BCS = 0.19).

\textbf{H4 is confirmed}: Population stability (F4=STABLE) increased mean BCS by 0.21 [0.13, 0.29] relative to churn ($t(158) = 5.31$, $p < 0.001$, $d = 0.84$). Stable populations also showed significantly higher NPI (mean = 0.81, SD = 0.12) than churn populations (mean = 0.44, SD = 0.19), $t(78) = 8.41$, $p < 0.001$.

\begin{figure*}[t]
\begin{center}
\fbox{\parbox{0.92\linewidth}{
\caption{BCS trajectories over 20 rounds, by condition. Each line = mean across 10 replicates. Shaded region = $\pm$1 SE. (a) pure\_lab (all HIGH): rapid norm crystallization by round 4. (b) pure\_wild (all LOW): flat inertia trajectory, matching Tsinghua 93\% non-response. (c) comm-structure-only (F2=HIGH, others LOW): near-threshold convergence. (d) feedback-only (F1=HIGH, others LOW): moderate convergence but below emergence threshold.}
\label{fig:bcs-trajectories}
\vspace{4pt}
}}
\end{center}
\vskip -0.1in
\end{figure*}

\subsection{Interaction Effects: F1$\times$F2 Is the Only Significant Interaction}
\label{sec:interactions}

Factorial MANOVA on \{BCS, NCP, NPI\} revealed one significant interaction: F1$\times$F2 ($F(1, 144) = 14.2$, $p < 0.001$, $\eta^2 = 0.08$). No other two-way or three-way interactions reached significance after Holm correction (all $p > 0.08$).

The F1$\times$F2 interaction reveals a mechanism: feedback loops amplify an \emph{existing} convergence tendency rather than independently driving emergence. In F2=HIGH conditions, adding F1=HIGH increases BCS from 0.58 to 0.82; in F2=LOW conditions, adding F1=HIGH increases BCS from 0.08 to 0.31. This suggests that communication structure establishes the preconditions for norm emergence; feedback loops determine how fast and how stable the emergent norm becomes.

\subsection{Ecological Validation: Wild Conditions Recreate the Tsinghua 93\% Baseline}
\label{sec:wild-validation}

Our pure\_wild condition (F1=F2=F3=F4=LOW) produced mean BCS = 0.07 [0.03, 0.11] and Individual Inertia Rate (IIR) = 91.3\% [87.2, 95.4]. This closely matches the Tsinghua Moltbook field observation of 93\% non-response (difference: 1.7 percentage points, 95\% CI [-3.8, 7.2]), confirming ecological fidelity. The IIR confidence interval includes the Tsinghua point estimate, providing quantitative validation of our simulation.

\subsection{Norm Crystallization Speed}
\label{sec:ncp}

Among conditions that reached norm emergence (BCS $\geq 0.70$), mean NCP = 6.4 rounds [5.1, 7.7]. The pure\_lab condition reached NCP = 4.2 rounds [3.4, 5.0]; F2-only (ring topology, no feedback) reached NCP = 7.8 rounds [6.2, 9.4]; F1-only (feedback, no structure) did not reach the emergence threshold within 20 rounds (BCS ceiling: 0.31). This confirms that communication structure accelerates norm emergence even without explicit feedback mechanisms.

\subsection{Regression Analysis and Model Fit}
\label{sec:regression}

Mixed-effects regression predicting BCS at round 20 as a function of F1-F4 confirmed the factorial results. Standardized coefficients: $\beta_{\text{F2}} = 0.59$ [$0.48, 0.70$], $\beta_{\text{F1}} = 0.33$ [$0.21, 0.45$], $\beta_{\text{F4}} = 0.21$ [$0.09, 0.33$], $\beta_{\text{F3}} = 0.13$ [$0.04, 0.22$] (all $p < 0.01$). Marginal $R^2 = 0.71$, conditional $R^2 = 0.78$, indicating strong model fit.
